\documentclass[11pt]{article}
\usepackage{raeez-math-template}
\title{Signed Archimedean Pairings and Canonical Determinants}
\author{Raeez Lorgat}
\date{}
\begin{document}
\maketitle

The logarithmic derivative of the gamma factor defines a signed form.
Its translation formula determines the Gaussian sign exactly.
A de Branges norm requires a separate map and a Hardy-space domain.
For a discrete operator, a canonical product fixes the divisor and two initial values.
These constructions give explicit pairing and determinant identities with their hypotheses stated.

\section{The signed translation form}

All vector spaces are over $\mathbb C$, and inner products are linear in the first variable.
Lebesgue measure on $\mathbb R$ is denoted by $dt$ or $du$.
For $f$ in the Schwartz space $\mathcal S(\mathbb R)$, use
\[
 h_f(u)=\mathcal F_+f(u)=\int_{\mathbb R}f(t)e^{iut}dt,
 \qquad
 f(t)=\frac1{2\pi}\int_{\mathbb R}h_f(u)e^{-iut}du.
\]
Define $f^*(t)=\overline{f(-t)}$, $\tau_t f(x)=f(x-t)$, and
$(f*g)(x)=\int f(v)g(x-v)dv$.
Then $\mathcal F_+(f*g^*)=h_f\overline{h_g}$ on the real line.
Put
\[
 w(u)=\log\pi-\Re\psi(1/4+iu/2),\qquad
 c_0=\log\pi-\psi(1/4),\qquad
 \kappa(t)=\frac{e^{-t/2}}{1-e^{-2t}}\quad(t>0),
\]
where $\psi=\Gamma'/\Gamma$ and $\gamma$ denotes Euler's constant.

\begin{theorem}[signed archimedean form]\label{thm:translation}
The continuous sesquilinear form
\begin{equation}\label{eq:pairing}
 B_\infty(f,g)=\frac1{2\pi}\int_{\mathbb R}
        h_f(u)\overline{h_g(u)}w(u)du
\end{equation}
on $\mathcal S(\mathbb R)$ is Hermitian. It satisfies
\begin{equation}\label{eq:translation}
 B_\infty(f,g)=c_0\langle f,g\rangle_{L^2}
 -\int_0^\infty\kappa(t)
       \langle\tau_t f-f,\tau_t g-g\rangle_{L^2}dt,
 \qquad
 c_0=\log\pi+\gamma+\frac\pi2+3\log2>0.
\end{equation}
Both integrals converge absolutely.
\end{theorem}
\begin{proof}
The digamma integral, valid for $\Re z>0$, is
\[
 \psi(z)+\gamma=\int_0^\infty
           \frac{e^{-v}-e^{-zv}}{1-e^{-v}}dv
 \quad\text{\cite[Eq.~5.9.16]{dlmf}}.
\]
Subtract its value at $z=1/4$, take real parts, and set $v=2t$.
This gives
\begin{equation}\label{eq:multiplier}
 w(u)=c_0-2\int_0^\infty\kappa(t)(1-\cos ut)dt.
\end{equation}
The inner integral is finite. Near zero, $1-\cos ut=O_u(t^2)$ and
$\kappa(t)=O(t^{-1})$. At infinity, $\kappa(t)=O(e^{-t/2})$.
The digamma series also gives
\begin{equation}\label{eq:positive-series}
 c_0-w(u)=\sum_{n=0}^\infty
 \frac{u^2/4}{(n+1/4)((n+1/4)^2+u^2/4)}.
\end{equation}
Each summand is nonnegative.
The series converges locally uniformly and is $O(1+\log(2+|u|))$.
For example, split it at $n=\lceil|u|\rceil$ and bound the remaining cubic tail.
Thus $w$ has logarithmic growth, so \eqref{eq:pairing} converges and is continuous on Schwartz space.
Its real multiplier makes it Hermitian.

Plancherel and Tonelli give \eqref{eq:translation} first for $g=f$.
The Fourier multiplier of $\tau_t-I$ has squared absolute value $2(1-\cos ut)$.
For general $f,g$, polarization gives the equality.
Absolute convergence also follows directly from
\[
 \|\tau_t f-f\|_2\le\min(t\|f'\|_2,2\|f\|_2).
\]
Finally, differentiation of gamma reflection and duplication gives
$\psi(1/4)=-\gamma-\pi/2-3\log2$.
These identities are \cite[Eqs.~5.5.4, 5.5.8]{dlmf}.
\end{proof}

For $\phi\in\mathcal S(\mathbb R)$, define
$W_\infty(\phi)=(2\pi)^{-1}\int(\mathcal F_+\phi)(u)w(u)du$.
Fourier inversion in the preceding proof gives two equivalent formulas:
\begin{align}
 W_\infty(\phi)
 &=c_0\phi(0)-\int_0^\infty\kappa(t)
                 (2\phi(0)-\phi(t)-\phi(-t))dt,\label{eq:distribution}\\
 &=(\log\pi+\gamma)\phi(0)
   +\int_0^\infty
    \frac{e^{-t/2}(\phi(t)+\phi(-t))-2e^{-2t}\phi(0)}{1-e^{-2t}}dt.
    \label{eq:subtraction}
\end{align}
For \eqref{eq:distribution}, the numerator vanishes to second order at zero.
For \eqref{eq:subtraction}, it vanishes to first order.
The difference between their constant terms equals the digamma integral at $1/4$.
In particular, if $\phi(0)=0$, then
$W_\infty(\phi)=\int_0^\infty\kappa(t)(\phi(t)+\phi(-t))dt$.
The subtraction at zero is part of the distribution.

\begin{proposition}[the density and the form have both signs]\label{prop:sign}
The even function $w$ is strictly decreasing on $(0,\infty)$.
It has exactly one positive zero.
The form $B_\infty$ takes both signs on nonzero compact smooth functions.
\end{proposition}
\begin{proof}
Differentiate \eqref{eq:positive-series} locally uniformly:
\[
 w'(u)=-\frac u2\sum_{n=0}^\infty
       \frac{n+1/4}{((n+1/4)^2+u^2/4)^2}<0\quad(u>0).
\]
The digamma asymptotic \cite[Eq.~5.11.2]{dlmf} gives
$w(u)=\log(2\pi)-\log|u|+O(|u|^{-2})$ as $|u|\to\infty$.
Since $w(0)=c_0>0$, continuity proves the first two assertions.
The Gaussian calculation below gives a Schwartz function with each sign.
If $\chi$ is a smooth cutoff equal to one near zero, then
$\chi(t/R)f(t)\to f(t)$ in Schwartz space for a Gaussian $f$.
Continuity of $B_\infty$ preserves either strict sign for all sufficiently large $R$.
\end{proof}

\section{Gaussian pairing and an explicit remainder}

For $\lambda>0$, put $f_\lambda(t)=e^{-\pi\lambda t^2}$.
Completing the square gives
\[
 h_{f_\lambda}(u)=\lambda^{-1/2}e^{-u^2/(4\pi\lambda)},\qquad
 (f_\lambda*f_\mu^*)(t)=\frac{e^{-\pi\lambda\mu t^2/(\lambda+\mu)}}{\sqrt{\lambda+\mu}}.
\]
In particular, $\|f_\lambda\|_2^2=(2\lambda)^{-1/2}$.
These functions belong to Schwartz space and do not have compact support.

\begin{theorem}[Gaussian sign threshold]\label{thm:gaussian}
Define
\begin{equation}\label{eq:gaussian}
 R(\lambda)=\frac{B_\infty(f_\lambda,f_\lambda)}{\|f_\lambda\|_2^2}
   =c_0-2\int_0^\infty\kappa(t)(1-e^{-\pi\lambda t^2/2})dt.
\end{equation}
Then $R$ is smooth and strictly decreasing on $(0,\infty)$, with
\[
 \lim_{\lambda\downarrow0}R(\lambda)=c_0,\qquad
 \lim_{\lambda\to\infty}R(\lambda)=-\infty.
\]
There is a unique $\lambda_c>0$ with $R(\lambda_c)=0$.
The Gaussian pairing is positive for $0<\lambda<\lambda_c$ and negative for $\lambda>\lambda_c$.
For $\lambda,\mu>0$, the polarized formula is
\begin{equation}\label{eq:gaussian-cross}
 B_\infty(f_\lambda,f_\mu)
   =\frac1{\sqrt{\lambda+\mu}}
          R\left(\frac{2\lambda\mu}{\lambda+\mu}\right).
\end{equation}
\end{theorem}
\begin{proof}
Substitute the Gaussian autocorrelation into \eqref{eq:distribution}.
Differentiation under the integral, locally uniformly in $\lambda>0$, gives
\[
 R'(\lambda)=-\pi\int_0^\infty t^2\kappa(t)e^{-\pi\lambda t^2/2}dt<0.
\]
All higher derivatives have integrable majorants on compact parameter intervals.
For the limit at zero, use
$1-e^{-\pi\lambda t^2/2}\le\pi\lambda t^2/2$ and $\int t^2\kappa(t)dt<\infty$.
For the limit at infinity, monotone convergence gives
\[
 \int_0^\infty\kappa(t)(1-e^{-\pi\lambda t^2/2})dt
       \longrightarrow\int_0^\infty\kappa(t)dt=\infty.
\]
The divergence occurs at zero, where $\kappa(t)\sim1/(2t)$.
The intermediate value theorem gives the unique threshold.
The convolution formula gives \eqref{eq:gaussian-cross}.
\end{proof}

\begin{proposition}[convergent complementary-error-function formula]\label{prop:erfc}
Put $\alpha=(2\pi\lambda)^{-1}$, $M=\sqrt{\pi/\alpha}$, and $c_n=n+1/4$.
Then
\begin{equation}\label{eq:erfc-series}
 R(\lambda)=c_0-\sum_{n=0}^\infty d_n(\lambda),\qquad
 d_n(\lambda)=\frac1{c_n}-\frac{2\pi}{M}
                  e^{4\alpha c_n^2}\operatorname{erfc}(2c_n\sqrt\alpha).
\end{equation}
Each $d_n$ is positive. For $N\ge1$,
\begin{equation}\label{eq:tail}
 0\le\sum_{n=N}^\infty d_n(\lambda)
    \le\frac{\pi\lambda}{8(N-1+1/4)^2}.
\end{equation}
\end{proposition}
\begin{proof}
For $a,\alpha>0$, the identity
$(a^2+u^2)^{-1}=\int_0^\infty e^{-v(a^2+u^2)}dv$ and Tonelli give
\[
 \int_{\mathbb R}\frac{e^{-\alpha u^2}}{a^2+u^2}du
 =\sqrt\pi\int_0^\infty\frac{e^{-a^2v}}{\sqrt{\alpha+v}}dv
 =\frac\pi a e^{\alpha a^2}\operatorname{erfc}(a\sqrt\alpha).
\]
The last step uses $r=\sqrt{\alpha+v}$.
Integrate \eqref{eq:positive-series} against $M^{-1}e^{-\alpha u^2}du$.
The normalized Gaussian second moment is $1/(2\alpha)=\pi\lambda$.
Tonelli gives \eqref{eq:erfc-series} and
$0<d_n\le\pi\lambda/(4c_n^3)$.
Finally,
$\sum_{n=N}^\infty c_n^{-3}\le\int_{N-1}^\infty(x+1/4)^{-3}dx$.
This proves \eqref{eq:tail}.
\end{proof}

The normalization in \eqref{eq:gaussian} can also be written as
\[
 R(\lambda)=\frac1{\sqrt{2\pi}}
          \int_{\mathbb R}e^{-v^2/2}w(\sqrt{\pi\lambda}\,v)dv.
\]
These signed formulas concern the archimedean term alone.
The full Weil form also contains its polar and finite-prime terms.

\section{A signed realization on the full zero multiset}

Let $\xi(s)=\tfrac12s(s-1)\pi^{-s/2}\Gamma(s/2)\zeta(s)$.
Let $\mathcal Z$ be its distinct zeros, with multiplicities $m_\rho$,
and put $z_\rho=-i(\rho-1/2)$.
Use $\mathscr A=C_c^\infty(\mathbb R)$ with its usual test-function topology.
Its Fourier transforms are entire.
Integration by parts gives rapid decay in every fixed horizontal strip.
The zero counting estimate $N(T)=O(T\log(2+T))$ therefore implies
\[
 \sum_{\rho\in\mathcal Z}m_\rho|h_f(z_\rho)|^2<\infty
 \qquad(f\in\mathscr A).
\]
The absolute zero sums and the Fourier convention are those of
\cite[Eqs.~(1.1)--(1.2)]{suzuki}.

\begin{proposition}[zero involution pairing]\label{prop:zero-involution}
On $\mathcal K=\ell^2(\mathcal Z,m)$, define
$(Jv)(\rho)=v(1-\bar\rho)$ and $(Sf)(\rho)=h_f(z_\rho)$.
Then $J$ is a self-adjoint unitary and
\begin{equation}\label{eq:zero-pairing}
 q_W(f,g):=\sum_{\rho\in\mathcal Z}m_\rho
 h_f(z_\rho)\overline{h_g(\bar z_\rho)}
       =\langle Sf,J Sg\rangle_{\mathcal K}.
\end{equation}
\end{proposition}
\begin{proof}
The functional equation and conjugation preserve the zero multiset.
The map $\rho\mapsto1-\bar\rho$ is an involution preserving multiplicity.
Moreover $z_{1-\bar\rho}=\bar z_\rho$.
Thus $J^2=I$, $J^*=J$, and direct substitution gives the identity.
Absolute convergence follows from Cauchy--Schwarz.
\end{proof}

If all zeros lie on the critical line, then $J=I$ and the zero form is nonnegative.
An off-line orbit gives a matrix $\left(\begin{smallmatrix}0&1\\1&0\end{smallmatrix}\right)$
on the corresponding coordinates of $\mathcal K$.
Its negative eigenspace alone does not identify the range $S(\mathscr A)$.
The converse positivity theorem requires Weil's criterion.

\section{The de Branges map and its domain}

Write $F^\sharp(z)=\overline{F(\bar z)}$.
Use the Hermite--Biehler condition $|E^\sharp(z)|<|E(z)|$ for $\Im z>0$.
This convention permits common real zeros of $E$ and $E^\sharp$.
The quotients below use their removable continuations.
Lagarias calls an $E$ without real zeros a strict de Branges function
\cite[Section~2, pp.~2--4]{lagarias}.
For such an $E$, put $\Theta=E^\sharp/E$ and
\[
 H^2=\left\{h\text{ holomorphic on }\mathbb C_+:
      \sup_{y>0}\int_{\mathbb R}|h(x+iy)|^2dx<\infty\right\},
 \qquad K_\Theta=H^2\ominus\Theta H^2.
\]
The de Branges space consists of entire $F$ with $F/E,F^\sharp/E\in H^2$,
and has norm $\|F\|_{H(E)}=\|F/E\|_{L^2}$.
Multiplication by $E$ is the isometry $K_\Theta\to H(E)$.
These conventions are given in \cite[Sections~2.3--2.4]{suzuki}.

\begin{example}[a zero of arbitrary multiplicity]\label{ex:multiple}
Let $A(z)=z^m$ for an integer $m\ge1$. Then
\[
 E=A+iA'=z^{m-1}(z+im),\qquad
 \Theta(z)=\frac{z-im}{z+im},\qquad
 e_0(z)=i\sqrt{\frac m\pi}\frac1{z+im}.
\]
The space $K_\Theta$ is one-dimensional, and $e_0$ is a unit vector because
$\int_{\mathbb R}(x^2+m^2)^{-1}dx=\pi/m$.
Also $e_0(0)=1/\sqrt{\pi m}$ and $\Theta'(0)=-2i/m$.
Thus $\|h\|_{K_\Theta}^2=\pi m|h(0)|^2$ on this space.
The multiplicity enters the sampling measure.
It is compatible with a single normalized basis vector at the distinct zero.
\end{example}

\begin{theorem}[explicit de Branges sampling realization]\label{thm:db}
Assume that all zeros of $\xi$ lie on the critical line.
Put $A(z)=\xi(1/2+iz)$, $E=A+iA'$, and $\Theta=E^\sharp/E$.
Let $\Gamma$ be the distinct real zeros of $A$, with multiplicities $m_\gamma$.
Then $E$ is Hermite--Biehler and
\begin{equation}\label{eq:basis}
 e_\gamma(z)=i\sqrt{\frac{m_\gamma}{\pi}}
       \frac{A(z)}{(z-\gamma)E(z)},\qquad \gamma\in\Gamma,
\end{equation}
form an orthonormal basis of $K_\Theta$.
For $f\in\mathscr A$, the series
\begin{equation}\label{eq:db-map}
 T f=\sum_{\gamma\in\Gamma}\sqrt{\pi m_\gamma}\,
              h_f(\gamma)e_\gamma\quad\text{in }K_\Theta,
 \qquad \mathcal L f=\pi^{-1/2}E\,Tf\quad\text{in }H(E)
\end{equation}
converges and satisfies
\begin{equation}\label{eq:db-pairing}
 \langle\mathcal L f,\mathcal L g\rangle_{H(E)}=q_W(f,g).
\end{equation}
For every $h\in K_\Theta$,
\begin{equation}\label{eq:sampling-norm}
 \|Eh\|_{H(E)}^2=\|h\|_{L^2}^2
            =\pi\sum_{\gamma\in\Gamma}m_\gamma|h(\gamma)|^2.
\end{equation}
\end{theorem}
\begin{proof}
The even order-one Hadamard product of $A$ gives, with positive zeros grouped,
\[
 A(z)=A(0)\prod_{\gamma>0}(1-z^2/\gamma^2)^{m_\gamma},\qquad
 L(z):=\frac{A'(z)}{A(z)}
   =\sum_{\gamma>0}m_\gamma
          \left(\frac1{z-\gamma}+\frac1{z+\gamma}\right).
\]
The grouped sum converges locally uniformly off the zeros because
$\sum m_\gamma/(1+\gamma^2)<\infty$.
The order and factorization are stated in \cite[Theorem~1 and proof, p.~9]{lagarias}.
Evenness removes the linear exponential in the Hadamard product.
There is no zero at the origin: $\zeta(1/2)\ne0$, as follows from
the positive alternating eta series and $1-2^{1/2}<0$.
For $\Im z>0$, each reciprocal has negative imaginary part.
Thus $\Im L(z)<0$ and
\[
 |E(z)|^2-|E^\sharp(z)|^2=-4|A(z)|^2\Im L(z)>0.
\]
This proves the Hermite--Biehler condition.

The boundary norm of $H^2$ uses $du$, so its kernel is
$[2\pi i(\bar w-z)]^{-1}$.
Orthogonal projection onto $K_\Theta$ therefore gives the model kernel
\[
 k_\Theta(w,z)=\frac{1-\overline{\Theta(w)}\Theta(z)}{2\pi i(\bar w-z)}.
\]
Since $\Theta=(1-iL)/(1+iL)$, direct algebra gives
\begin{align*}
 k_\Theta(w,z)
 &=\frac{L(z)-\overline{L(w)}}{
       \pi(\bar w-z)(1+iL(z))(1-i\overline{L(w)})}\\
 &=\sum_{\gamma\in\Gamma}e_\gamma(z)\overline{e_\gamma(w)}.
\end{align*}
The difference of each pair of reciprocal terms cancels its linear denominator.
The resulting series is absolutely locally convergent for $z,w\in\mathbb C_+$.

Near $\gamma$, write $A(z)=(z-\gamma)^{m_\gamma}a(z)$ with $a(\gamma)\ne0$.
Then
\[
 \frac{A(z)}{E(z)}=\frac{z-\gamma}{im_\gamma}+O((z-\gamma)^2),\qquad
 \Theta'(\gamma)=-\frac{2i}{m_\gamma}.
\]
It follows that
$e_\gamma(\gamma)=1/\sqrt{\pi m_\gamma}$ and
$e_\gamma(\delta)=0$ for every distinct zero $\delta$.
The kernel has the finite boundary limit
$k_\Theta(\gamma,z)=e_\gamma(z)/\sqrt{\pi m_\gamma}$.
These boundary evaluations are bounded: the kernels at $\gamma+iy$
converge in norm, as the inner products of these kernels converge to their finite analytic limits.
Consequently the $e_\gamma$ are normalized reproducing kernels and are orthonormal.

For fixed $w\in\mathbb C_+$, the kernel expansion on its diagonal gives
$\sum_\gamma|e_\gamma(w)|^2=k_\Theta(w,w)$.
The orthogonal projection of $k_\Theta(w,\cdot)$ onto their closed span
therefore has the full norm of that kernel.
All reproducing kernels lie in the span, proving completeness.
Parseval now gives \eqref{eq:sampling-norm}.
The rapid strip decay of $h_f$ gives
$\sum m_\gamma|h_f(\gamma)|^2<\infty$, so \eqref{eq:db-map} converges in norm.
Its coefficients give
$\langle Tf,Tg\rangle=\pi q_W(f,g)$.
Multiplication by $E$ and division by $\sqrt\pi$ prove \eqref{eq:db-pairing}.
\end{proof}

The basis and its normalization also appear in
\cite[Eq.~(3.5), Proposition~4.1]{suzuki}.
The theorem assumes real zeros before constructing a positive realization.
Conversely, if this $E$ is Hermite--Biehler, a nonreal zero of $A$ is impossible.
At an upper-half-plane zero, $E=iA'$ and $E^\sharp=-iA'$ have equal absolute value.
Real conjugation excludes lower-half-plane zeros as well.
Thus this Hermite--Biehler condition contains the real-zero requirement.

\begin{proposition}[sampling mass and spectral multiplicity]\label{prop:multiplicity}
Under the hypotheses of Theorem~\ref{thm:db}, the diagonal operator
\[
 A_K\Bigl(\sum_\gamma a_\gamma e_\gamma\Bigr)
       =\sum_\gamma\gamma a_\gamma e_\gamma,\qquad
 D(A_K)=\left\{\sum_\gamma a_\gamma e_\gamma:
                \sum_\gamma\gamma^2|a_\gamma|^2<\infty\right\}
\]
is self-adjoint and has simple eigenvalues at the distinct zeros $\gamma$.
The sampling factor $m_\gamma$ does not change that multiplicity.
\end{proposition}
\begin{proof}
The orthonormal basis identifies $A_K$ with real diagonal multiplication on $\ell^2(\Gamma)$.
Its eigenspace at $\gamma$ is the one-dimensional span of $e_\gamma$.
Thus its resolvent determinant has one zero at $1/2+i\gamma$.
If $m_\gamma>1$, the full zero divisor requires additional eigendirections.
An explicit enlargement is $\bigoplus_\gamma\mathbb C^{m_\gamma}$,
with diagonal eigenvalue $\gamma$ on each summand.
The sampling vector has equal coordinates $h_f(\gamma)$ within that summand.
Its squared norm is $m_\gamma|h_f(\gamma)|^2$.
The orthogonal complement of the equal-coordinate line has dimension $m_\gamma-1$.
It contributes to the determinant while remaining orthogonal to every sampling vector.
\end{proof}

\begin{proposition}[the Gaussian domain obstruction]\label{prop:hardy}
For every $\lambda>0$, $h_{f_\lambda}\notin H^2(\mathbb C_+)$.
Consequently $E h_{f_\lambda}\notin H(E)$ for every Hermite--Biehler $E$.
\end{proposition}
\begin{proof}
The Gaussian Fourier formula gives
\[
 \int_{\mathbb R}|h_{f_\lambda}(x+iy)|^2dx
    =\pi\sqrt{\frac2\lambda}\,e^{y^2/(2\pi\lambda)}.
\]
This is unbounded as $y\to\infty$.
Membership of $E h_{f_\lambda}$ in $H(E)$ would require its quotient by $E$ to belong to $H^2$.
\end{proof}

For $h\in K_\Theta$, the map $h\mapsto Eh$ transports the Lebesgue norm exactly.
Transport of $B_\infty$ instead gives its weighted multiplier $w/(2\pi)$ on the domain
$\int(1+|w|)|h|^2<\infty$.
Equation~\eqref{eq:pairing} and Proposition~\ref{prop:sign} show why that signed form cannot equal a positive norm on all compact tests.

\section{Canonical genus-one determinant}

Let $H$ be a separable complex Hilbert space.
Assume a specified self-adjoint operator $A$ has an orthonormal eigenbasis $(v_j)$,
real eigenvalues $\lambda_j$, and finite multiplicities.
Its domain is
\[
 D(A)=\left\{\sum_j a_jv_j:\sum_j\lambda_j^2|a_j|^2<\infty\right\}.
\]
Assume $\sum_j(1+\lambda_j^2)^{-1}<\infty$.
Define $\Theta=\tfrac12 I+iA$ on $D(A)$, $\rho_j=\tfrac12+i\lambda_j$, and
$T=\Theta^{-1}$.
Then $T$ is bounded and Hilbert--Schmidt.
For a Hilbert--Schmidt operator $K$, use
$\det_2(I+K)=\det((I+K)e^{-K})$.
The operator inside the last determinant differs from the identity by a trace-class operator.

\begin{theorem}[product and regularized resolvent]\label{thm:det}
The function
\begin{equation}\label{eq:det}
 D(s)=\det_2(I-sT)=\prod_j(1-s/\rho_j)e^{s/\rho_j}
\end{equation}
is entire, with its zeros exactly the $\rho_j$, counted with multiplicity.
It satisfies $D(0)=1$, $D'(0)=0$, and, off the zero set,
\begin{align}
 \frac{D'(s)}{D(s)}
   &=\operatorname{Tr}\bigl((sI-\Theta)^{-1}+\Theta^{-1}\bigr),\label{eq:resolvent}\\
 \left(\frac{D'}D\right)'(s)
   &=-\operatorname{Tr}\bigl((sI-\Theta)^{-2}\bigr).\label{eq:second}
\end{align}
The sum in \eqref{eq:resolvent} is trace class as one operator.
\end{theorem}
\begin{proof}
Write $E_1(z)=(1-z)e^z$. For $|z|\le1/2$,
$\log E_1(z)=-\sum_{k\ge2}z^k/k$ and $|\log E_1(z)|\le|z|^2$.
On each compact set of $s$, the tails of \eqref{eq:det} therefore converge uniformly.
Finite initial factors show the claimed divisor.
Diagonalizing $(I-sT)e^{sT}$ identifies its Fredholm determinant with the product.
At zero, every logarithmic linear term vanishes.
Termwise logarithmic differentiation gives
\[
 \frac{D'}D=\sum_j\left(\frac1{s-\rho_j}+\frac1{\rho_j}\right),\qquad
 \left(\frac{D'}D\right)'=-\sum_j\frac1{(s-\rho_j)^2}.
\]
Both series converge absolutely and locally uniformly off the divisor.
Indeed their tails are bounded by constant multiples of $\sum_j|\rho_j|^{-2}$.
The resolvent identity gives
\[
 (sI-\Theta)^{-1}+\Theta^{-1}
      =s(sI-\Theta)^{-1}\Theta^{-1}.
\]
Both factors on the right are Hilbert--Schmidt, so their product is trace class.
Taking traces in the eigenbasis proves the identities.
\end{proof}

\begin{proposition}[base-point normalization]\label{prop:basepoint}
Suppose the multiset $(\rho_j)$ equals the zero divisor of $\xi$.
Then
\begin{equation}\label{eq:xi-det}
 \xi(s)=\frac12e^{bs}\det_2(I-s\Theta^{-1}),\qquad
 b=\log(2\sqrt\pi)-1-\frac\gamma2.
\end{equation}
More generally, for $s_0$ off the divisor,
\begin{equation}\label{eq:basepoint}
 \frac{\xi(s)}{\xi(s_0)}
  =\exp\left((s-s_0)\frac{\xi'(s_0)}{\xi(s_0)}\right)
       \det_2\left(I+(s-s_0)(s_0I-\Theta)^{-1}\right).
\end{equation}
\end{proposition}
\begin{proof}
The order-one Hadamard factorization of $\xi$ gives
$\xi(s)=e^{a+bs}\prod_j E_1(s/\rho_j)$.
See \cite[proof of Theorem~1, p.~9]{lagarias}.
At zero, $e^a=\xi(0)$ and $b=\xi'(0)/\xi(0)$.
The expansions
\[
 \Gamma(s/2)=2/s-\gamma+O(s),\qquad
 \zeta(s)=-\frac12-\frac12\log(2\pi)s+O(s^2)
\]
give $\xi(0)=1/2$ and the displayed $b$.
The zeta values are \cite[Eqs.~25.6.1, 25.6.11]{dlmf}.
For the base-point formula, both logarithmic second derivatives equal
$-\sum_j(s-\rho_j)^{-2}$.
The ratio of the two sides is holomorphic and nonvanishing across their identical zeros.
Its logarithmic derivative is constant on the connected complement of the divisor.
The value and first logarithmic derivative at $s=s_0$ make that ratio one.
\end{proof}

The eigenvalue hypothesis in Proposition~\ref{prop:basepoint} is separate from the normalization.
It supplies the full divisor before \eqref{eq:xi-det} is deduced.
The two initial values fix the remaining exponential factor.
The determinant formula does not construct an operator with the required eigenvalues.

\begin{example}[spectral zeta and canonical normalization]
On a one-dimensional space, let $\Theta=\rho I$ with $\rho\ne0$.
On any domain with a chosen branch of $\log(s-\rho)$, the spectral zeta function is
$Z(u;s)=(s-\rho)^{-u}$, so
\[
 e^{-\partial_u Z(0;s)}=s-\rho,
 \qquad
 \det_2(I-s\Theta^{-1})=(1-s/\rho)e^{s/\rho}.
\]
Their ratio is $-\rho e^{-s/\rho}$.
The divisor alone does not identify these two normalizations.
\end{example}

A spectral zeta determinant requires a spectral cut, a domain for its powers,
and meromorphic continuation regular at the differentiation point.
For self-adjoint operators bounded below, \cite[Theorem~4.5, Eq.~(4.7)]{hartmann-lesch}
states the additional analytic hypotheses and the exponential polynomial relating the determinants.
The normal operator $\Theta=\tfrac12+iA$ requires its own such analysis.
The canonical product in Theorem~\ref{thm:det} already gives
\eqref{eq:xi-det} under the stated divisor hypothesis.

\begin{thebibliography}{9}
\bibitem{dlmf}
NIST Digital Library of Mathematical Functions.
\href{https://dlmf.nist.gov/5.7.E6}{Equation 5.7.6},
\href{https://dlmf.nist.gov/5.9.E16}{Equation 5.9.16},
\href{https://dlmf.nist.gov/5.11.E2}{Equation 5.11.2}, and
\href{https://dlmf.nist.gov/5.5}{Section 5.5}.
Zeta values: \href{https://dlmf.nist.gov/25.6.E1}{Equation 25.6.1} and
\href{https://dlmf.nist.gov/25.6.E11}{Equation 25.6.11}.

\bibitem{suzuki}
M.~Suzuki, \emph{On the Hilbert space derived from the Weil distribution}.
\href{https://arxiv.org/html/2301.00421v3}{arXiv:2301.00421v3}.
Sections 2.3--2.4, equation (3.5), Proposition 4.1, and Theorem 1.1.

\bibitem{hartmann-lesch}
L.~Hartmann and M.~Lesch,
\emph{Zeta and Fredholm determinants of self-adjoint operators}.
\href{https://arxiv.org/html/2106.02444v2}{arXiv:2106.02444v2}.
Theorems 4.3 and 4.5, equations (4.4)--(4.7).

\bibitem{lagarias}
J.~C.~Lagarias,
\emph{Hilbert Spaces of Entire Functions and Dirichlet $L$-Functions}.
In \emph{Frontiers in Number Theory, Physics and Geometry I},
Springer, 2006, pp.~365--377.
\href{https://websites.umich.edu/~lagarias/doc/debranges-houches.pdf}{Author's version},
Section~2, pp.~2--4 and 7, and Theorem~1, p.~9.
\end{thebibliography}
\end{document}
